\section{Projekt trasowy i lokalizacja urządzeń}
\subsection{Plan strefy terminalowo-technicznej}
\fig{site-plan.pdf}{1}{Plan modelowego obszaru 600 $\times$ 400 m. Układ lokalny w metrach, północ umowna. Obrysy budynków odwzorowują przyjęte wymiary; oznaczenia urządzeń pokazują zespoły lokalne.}
\begin{tabularx}{\linewidth}{@{}P{25mm}P{38mm}Y@{}}
\toprule\textbf{Węzeł / (x,y)} & \textbf{Kamery} & \textbf{Pozostałe urządzenia}\\\midrule
N0 / (280,120) & C01--C08 wewn.; C09--C10 zewn. & SOC, rdzeń, VMS, AI, CSP-0; D01--D06; SSWiN-0.\\
N1 / (100,260) & C11--C14 stałe, C15 PTZ & CSP-1, D07--D09, SSWiN-1, szafa i UPS.\\
N2 / (470,240) & C16--C18 stałe, C19 PTZ & CSP-2, D10--D11, SSWiN-2, agregat w wydzielonym miejscu obok budynku.\\
N3 / (520,40) & C20--C22 stałe, C23--C24 PTZ & CSP-3, D12, G01--G02, SSWiN-3, szafa i UPS.\\\bottomrule
\end{tabularx}
Zespoły kamer zewnętrznych obejmują elewacje, bramy i wybrane kierunki obserwacji, a nie całą powierzchnię terenu. Dokładny punkt montażu, wysokość i kierunek kadru ustala się po próbie terenowej. Kamery nie są połączone długimi odcinkami miedzi do odległej szafy; w razie przekroczenia 90 m potrzebny jest dodatkowy punkt dystrybucyjny i aktualizacja materiałów.
